\chapter{Einleitung}
\label{chap:einleitung}

Welche Argumente können gemeinsam vertreten werden, wenn die zugrunde liegenden Informationen einander widersprechen? 
Diese Frage ist für menschliche ebenso wie für maschinelle Entscheidungsprozesse von zentraler Bedeutung. 
Abstrakte Argumentationssysteme (engl.\ \emph{abstract argumentation frameworks}, kurz AFs) stellen hierfür ein einfaches und zugleich allgemeines Modell bereit. 
Sie wurden von \textcite{dung:1995:af} eingeführt und bilden Argumente als Knoten sowie Angriffe zwischen ihnen als gerichtete Kanten ab. 

Ein Konflikt lässt jedoch nicht notwendigerweise nur eine einzige vertretbare Bewertung zu. 
Für dasselbe Argumentationssystem können unterschiedliche zulässige Bewertungen der enthaltenen Argumente bestehen. 

AN DIESER STELLE BEREITS SEMANTIKEN MIT EINEM BEISPIEL EINFÜHREN

Genau an dieser Stelle wird die bedingte Unabhängigkeit relevant. Sie stammt ursprünglich aus der Wahrscheinlichkeitstheorie und beschreibt, 
ob zwei Gruppen von Variablen noch Informationen übereinander liefern, nachdem eine dritte Gruppe beobachtet wurde. 
Grundlegende Eigenschaften solcher Unabhängigkeitsrelationen wurden unter anderem von \textcite{pearl:paz:1986:graphoids} untersucht. 
Eine bekannte Anwendung bilden Bayes-Netze, in denen Unabhängigkeiten zwischen Variablen mithilfe einer graphischen Darstellung 
sichtbar gemacht werden können und für die strukturierte Verarbeitung probabilistischer Informationen genutzt werden \cite{pearl:2010}.

Die hinter diesem Begriff stehende Idee ist jedoch nicht auf Wahrscheinlichkeiten beschränkt. Auch logische Wissensrepräsentationen 
können mehrere mögliche Zustände oder Interpretationen zulassen. Die von \textcite{darwiche:1997:logical-ci} entwickelte logische 
bedingte Unabhängigkeit überträgt die zentrale Frage daher auf das logische Schließen: KONKRETER, CI INFORMELL DEFINIEREN
Können Informationen aus unterschiedlichen 
Bereichen miteinander kombiniert werden, sobald eine gemeinsame Hintergrundinformation festgelegt ist?

In abstrakten Argumentationssystemen treten an die Stelle möglicher Interpretationen die zulässigen Bewertungen der Argumente.
Zwei Argumentmengen können als bedingt unabhängig angesehen werden, wenn ihre möglichen Bewertungen nach Festlegung einer 
dritten Argumentmenge miteinander vereinbar bleiben. Bedingte Unabhängigkeit erfasst damit, ob zwei argumentative 
Teilbereiche getrennt voneinander betrachtet werden können.

Eine entsprechende Definition bedingter Unabhängigkeit für abstrakte Argumentationssysteme wurde von \citeauthor*{rienstra:2020:independence-af} \cite{rienstra:2020:independence-af} eingeführt. 
Das damit verbundene Potenzial liegt auf der Hand: Werden unabhängige Teilbereiche erkannt, können argumentative Zusammenhänge verständlicher dargestellt 
und umfangreiche Probleme in kleinere Teilprobleme zerlegt werden. Diesem Potenzial steht jedoch eine wesentliche Schwierigkeit gegenüber. 
Die exakte Entscheidung bedingter Unabhängigkeit ist für zentrale Auswertungsweisen abstrakter Argumentationssysteme bereits mindestens auf der 
zweiten Ebene der polynomiellen Hierarchie anzuordnen \cite{bluemel:2025:independence-aba}.

HIER ERWÄHNEN, WAS AUSWERTUNG IST UND, DASS DESHALB CI SO SCHWER IST. 
Für die automatisierte Auswertung formaler Argumentationsprobleme wird daher häufig der Reduktionsansatz verwendet. 
Statt eine Aufgabe unmittelbar zu lösen, wird sie in einen anderen Formalismus übersetzt, für den bereits leistungsfähige 
allgemeine Solver verfügbar sind. Aussagenlogische Kodierungen können beispielsweise von SAT-Solvern verarbeitet werden. 
Quantifizierte boolesche Formeln (QBF) erweitern diesen Ansatz um die Eigenschaft Beziehungen zwischen mehreren möglichen Bewertungen 
innerhalb einer einzigen Formel auszudrücken, und eignen sich daher besonders für die Formulierung bedingter Unabhängigkeiten.

Dass QBFs als einheitliches Zielformat für anspruchsvolle Aufgaben der formalen Argumentation eingesetzt werden können, 
zeigen bereits \textcite{egly:woltran:2006:reasoning-qbf}.
Mit diesem Formalismus gewappnet wollen wir uns nun der Untersuchung der bedingten Unabhängikeit in Argumentationssystemen stellen.

\section{Zielsetzung}
\label{sec:einleitung-motivation}

Ziel dieser Arbeit ist es, bedingte Unabhängigkeit in abstrakten Argumentationssystemen sowohl praktisch als auch theoretisch zu untersuchen. 
Dazu wird die Unabhängigkeitsentscheidung für die \emph{stable}-, \emph{complete}- und \emph{preferred}-Semantik als QBF formuliert.
Die Kodierung ermöglicht eine exakte und automatisierte Prüfung, ohne die zulässigen Bewertungen eines Argumentationssystems vorab vollständig aufzählen zu müssen. 
Dadurch wird eine empirische Analyse auf größeren Benchmarkinstanzen ermöglicht.

Von besonderem Interesse ist die Konditionierungsmenge \(Z\); sie bestimmt, welche Argumentbewertungen bereits als festgelegt gelten. 
Eine Veränderung ihrer Größe kann beeinflussen, ob zwei Argumentmengen unabhängig sind und wie aufwendig die entsprechende Solveranfrage ist. 
Die Untersuchung auf ICCMA-Instanzen soll zeigen, ob sich hierbei systematische Zusammenhänge beobachten lassen.

Ergänzend wird analysiert, welchen Einfluss die Graphstruktur eines Argumentationssystems auf die Schwierigkeit der Unabhängigkeitsentscheidung besitzt. 
Strukturelle Einschränkungen können die Auswertung eines Argumentationssystems vereinfachen. Dies muss jedoch nicht unmittelbar für bedingte Unabhängigkeit gelten, 
da diese die Beziehungen zwischen mehreren zulässigen Bewertungen betrachtet. Die komplexitätstheoretische Analyse soll daher feststellen,
in welchen Graphklassen tatsächlich eine Vereinfachung erreicht wird und in welchen die Schwierigkeit des allgemeinen Falls bestehen bleibt.

\section{Beitrag der Arbeit}
\label{sec:einleitung-roadmap}
HIER BEREITS ERFOLGE UND ERGEBNISSE DER ARBEIT AUFZÄHLEN
\begin{itemize}
    \item Kapitel~2 führt die formalen Grundlagen der Arbeit ein. Dazu gehören abstrakte Argumentationssysteme, die betrachteten Auswertungsweisen, 
    bedingte Unabhängigkeit, quantifizierte boolesche Formeln sowie die benötigten komplexitätstheoretischen Begriffe.
    
    \item
    Kapitel~3 entwickelt labellingbasierte Kodierungen der \emph{stable}-, \emph{complete}- und \emph{preferred}-Semantik.
    Diese sind Ausgangspunkt für die anschließende QBF-Kodierung der bedingten Unabhängigkeit unter der \emph{stable}-, \emph{complete}- und \emph{preferred}-Semantik.
    

    \item
    Kapitel~4 Nutzt die zuvor entwickelten QBF-Kodierungen der bedingten Unabhängigkeit, um den Zusammenhang zwischen der Konditionierungsmenge und dem Unabhängigkeitsstatu
    beziehungsweise der Solverlaufzeit zu analysieren. 

    \item
    Kapitel~5 analysiert die Komplexität der bedingten Unabhängigkeit für ausgewählte Graphklassen. Wir zeigen, dass bedingte Unabhängigkeit für azyklische und geradezykelfreie
    AFs allgemeingültig ist und für ungeradezykelfreie und bipartite AFs die Schwierigkeit des allgemeinen Falls behält.
\end{itemize}

